\documentclass[UTF8,a4paper,12pt]{ctexart}
\usepackage{geometry}
\usepackage{hyperref}
\usepackage{url}
\usepackage{enumitem}
\usepackage{setspace}
\usepackage{titlesec}
\usepackage{indentfirst}
\usepackage{amssymb}
\usepackage{newunicodechar}
\newunicodechar{ψ}{$\psi$}
\newunicodechar{□}{$\square$}
\geometry{left=2.8cm,right=2.8cm,top=2.6cm,bottom=2.6cm}
\hypersetup{colorlinks=true,linkcolor=black,citecolor=black,urlcolor=blue}
\setstretch{1.5}
\setlength{\parindent}{2em}
\setlist{nosep}
\title{基于人工智能的天线优化设计研究综述}
\author{}
\date{\today}
\begin{document}
\maketitle
\tableofcontents
\newpage
\begin{abstract}

随着无线通信与雷达探测技术的飞速发展，现代工程对天线的综合性能提出了更高要求。传统基于数值全波电磁仿真的天线设计方法面临计算成本高、寻优周期长以及易陷入局部最优等瓶颈。为突破上述局限，探讨人工智能（AI）在天线设计领域的应用潜力与演进趋势具有重要意义。本文对AI辅助天线设计的相关研究进行了全面综述。首先，系统梳理了人工神经网络（ANN）、深度神经网络（DNN）以及图神经网络（GNN）在天线性能预测、复杂结构表征与阵列综合中的应用现状；其次，探讨了元启发式算法与深度学习代理模型协同优化技术的发展历程。研究发现，AI技术显著提升了复杂天线设计的非线性建模能力与多目标优化效率，但在高质量数据依赖、跨结构泛化能力、物理可解释性及工程可制造性验证方面仍存在明显局限。同时，本综述指出了当前领域在标准化评价体系、物理机理深度融合及多物理场协同等方面的研究空白。基于上述分析，本文提出未来AI天线设计应彻底告别碎片化的“一任务一模型”专用范式，加速向着以大规模电磁数据预训练为核心的“天线基础模型”演进，以实现跨结构与跨频段的知识迁移；同时亟需推动图深度学习向“全域电磁拓扑”延伸，并依托大语言模型智能体与主动学习机制构建打通“需求-生成-仿真-反馈”的智能闭环设计平台。本综述为全面实现天线设计的通用化、全域化与全流程自动化落地提供了理论参考与创新思路。

\par\noindent\textbf{关键词：}AI辅助设计、深度学习代理模型、图神经网络、天线基础模型、多目标优化

\end{abstract}

\section{引言}
随着无线通信、雷达探测、物联网以及 5G/6G 通信技术的高速发展，现代工程应用对天线综合性能提出了愈发严苛的要求，多频段、超宽带、高增益、低旁瓣、小型化等指标已成为天线设计中的主流需求。现阶段，HFSS、CST 等基于数值分析方法的电磁仿真软件被广泛应用于天线设计与优化过程，是开展电磁分析不可或缺的重要工具。然而，此类仿真软件在实际应用中存在明显局限：一方面，全波电磁求解计算量庞大，对硬件资源要求极高；另一方面，高度依赖工程师经验的“试错式”参数调试繁琐且重复性强，在面对多结构参数的批量寻优时，极易陷入局部最优，综合性能往往难以达到设计预期[1-3]。

为缓解传统设计中计算成本高、迭代周期长的问题，以人工神经网络（Artificial Neural Network, ANN）和深度神经网络（Deep Neural Network, DNN）为代表的深度学习技术逐步被引入天线设计领域。早在 1998 年，Mishra 等人便率先验证了神经网络在贴片天线性能预测中的可行性[14]。发展至今，深度学习凭借其优异的非线性拟合与特征表征能力，能够精准挖掘天线“结构—性能”间的内在映射关系。通过构建轻量化的代理模型（Surrogate Model），该类方法能够在极短时间内实现对未知天线结构电磁响应的快速预测，大幅缩短性能评估周期[4-9]。

然而，仅依靠代理模型实现“正向预测”尚不足以完成天线的“逆向设计”。在庞大且非凸的设计空间中，如何依据既定指标反向推演最优结构参数，是天线智能优化的核心环节。为此，遗传算法（Genetic Algorithm, GA）、粒子群优化（Particle Swarm Optimization, PSO）以及各类新型群智能算法等元启发式优化算法（Meta-heuristic Algorithms）发挥了关键作用。这类算法不依赖明确的梯度信息，擅长在复杂约束和多目标场景下开展全局搜索。近年来，业界已逐渐演变出“代理模型快速评估 + 元启发式全局寻优”的协同设计框架。在该框架下，代理模型替代高耗时的全波仿真承担性能预测，元启发式算法则驱动候选结构的迭代更新。这种结合有效兼顾了计算效率与全局寻优能力，成为当前天线设计优化的重要技术路线[10-13]。

尽管上述协同框架在单一天线或规则阵列设计中取得了显著成效，但在面对非均匀稀疏阵列、复杂共形阵列等具有多样化拓扑结构的现代天线系统时，传统模型仍显现出局限性。传统 ANN 和 DNN 通常只能处理固定维度的欧式空间张量数据，难以准确表征阵列中非规则的空间排布与复杂的电磁互耦合效应。近年来，图神经网络（Graph Neural Network, GNN）凭借对非欧式空间图结构数据的强大建模能力，为这一难题提供了全新解法。GNN 能够将天线阵元、物理间距、馈电与耦合关系直接抽象为图的节点与拓扑边，突破固定拓扑限制，实现对复杂阵列拓扑关系的高保真建模[19-23]。

由此可见，现代 AI 辅助天线设计已不再是单一算法的孤立应用，而是由深度学习特征提取、元启发式全局寻优以及图神经网络拓扑建模共同推动的设计范式转变。探索如何将几何结构、电磁响应、多目标约束与智能搜索机制统一于更加高效的通用表征与闭环框架内，进一步提升模型在复杂天线系统中的泛化能力，已成为该领域未来发展的重要方向。

\section{文献方法综述}

\subsection{基于深度学习的天线设计研究现状}

传统的天线设计和优化多依赖于工程师的理论基础和工程经验，通常需要使用电磁仿真软件通过等步长扫参方式进行多次调试，优化结果难以保证全局最优，并且设计过程复杂、耗时较长。随着人工智能技术的迅速发展，深度学习方法逐渐被引入天线领域。该类方法通过学习天线结构参数与性能指标之间的映射关系，在一定程度上弥补了传统电磁仿真软件计算成本高、设计周期长等不足，为天线设计与优化提供了新的技术路径。

目前，人工神经网络（Artificial Neural Network, ANN）和深度神经网络（Deep Neural Network, DNN）是在天线领域应用较多且较为成熟的网络模型，主要研究方向包括天线结构设计、性能预测以及天线阵列综合等。在天线结构设计和优化方面，相关研究多采用 ANN 进行训练和预测。例如，Mishra R K 等人较早将人工神经网络应用于方形贴片天线设计，验证了人工智能方法用于天线设计优化的可行性[14]。随后，为提高模型精度和训练效率，许多研究在基础 ANN 模型上进行改进，构建了更适用于天线设计任务的网络模型。

在具体应用方面，近年来的研究更加关注低样本代理建模、复杂结构快速生成和多目标自动设计。例如，Mahouti 等人针对透射阵天线提出低成本代理建模方法，利用自动化深度学习构建单元模型，在较少训练样本条件下完成 8-14 GHz、22-28 GHz 和 28-36 GHz 三个频段透射阵设计，并通过实验进行了验证[15]。Wei 等人提出 CNN-LSTM 结合的天线结构自动建模方法，仅以论文或图像中的天线结构图作为输入，便可自动生成对应的建模代码，从而加快不同天线物理结构样本的数据获取过程[16]。另有研究面向任意形状多边形贴片天线，采用 CNN 代理模型结合多岛差分进化算法进行多目标逆向设计，在数分钟内完成对反射系数、输入阻抗和辐射方向图等多个目标的综合优化[10]。

除单一天线结构设计外，深度学习方法也被应用于天线阵列的设计和综合。Ayestaran R G 等人使用神经网络进行非均匀天线阵列合成，将天线阵列的辐射模式与各阵元激励电压相关联，在不对阵列几何形状作过多假设的情况下实现阵列综合，并可在考虑耦合效应时处理非均匀阵列问题[17]。Mishra Subhash 等人利用 RBF 神经网络对共线短偶极子和平行短偶极子的均匀线性阵列进行方向性估计，该方法能够适应天线阵列的非线性特性，计算过程相对简洁，具有较好的鲁棒性[18]。

随着网络结构的不断发展，深度学习代理模型不再局限于简单的 S 参数预测，而是逐渐与多目标优化、结构生成和数据自动获取结合。相关研究表明，深度网络可作为快速预测器嵌入进化算法中，用于处理具有大量设计变量和多个性能指标的现代天线综合问题；同时，CNN、LSTM 等网络也可服务于训练样本自动生成，缓解传统手工建模效率较低的问题。这说明深度学习在天线设计中的作用正在从“性能预测工具”扩展为“建模、预测和优化协同工具”。

总体来看，深度学习方法在天线设计与优化中具有较强的非线性建模能力，能够在一定程度上替代重复的电磁仿真计算，提高天线结构设计、性能预测和阵列综合效率。然而，传统 ANN 和 DNN 主要适用于规则欧式数据或固定拓扑结构，对于非均匀阵列、复杂拓扑阵列等非欧式结构数据的表达能力仍存在一定局限，这也推动了图深度学习方法在天线设计领域中的进一步应用。

\subsection{基于图深度学习的天线设计研究现状}

为更好地学习非欧式空间数据的特征，利用深度学习技术实现图数据的端到端建模已成为重要研究方向。图神经网络（Graph Neural Network, GNN）能够处理由节点和边构成的图结构数据，适合描述复杂对象之间的连接关系和相互作用。随着图卷积、图注意力等方法的发展，越来越多基于图的神经网络被应用于社交网络、交通预测、生物制药以及电路和电磁建模等领域。

在天线和阵列问题中，阵元、用户、馈电点或可移动天线位置之间往往存在天然的图关系，因此可用图神经网络描述节点之间的连接、距离和耦合关系。近年来，GNN 已从早期的电路和阵列性能预测逐渐扩展到稀疏阵列测向、流体天线系统和新型可重构天线系统的联合优化中。

例如，Yang 等人针对随机稀疏线阵的波达方向估计问题，提出基于 GraphSAGE 的残差图卷积网络，将稀疏阵列接收信号构造为图结构，通过邻居节点聚合与更新补偿缺失阵元信息。该方法在低信噪比、低快拍数和大稀疏度条件下仍表现出较好的鲁棒性，在特定条件下相较 MUSIC 算法取得了明显性能提升[19]。

此外，GNN 也开始用于新型可重构天线系统的联合设计。He 等人面向流体天线系统提出两阶段 GNN 框架，分别推断天线位置和波束形成向量，并通过无监督损失函数联合训练，以实现多用户系统中的和速率与能效优化[22]。Xie 等人则针对 Pinching 天线系统，将下行通信系统建模为二部图，提出基于图注意力网络的 BGAT 模型，用于联合优化天线放置和功率分配，并强调了模型在最优性、可扩展性和计算效率方面的优势[23]。

综上，图神经网络为天线阵列建模提供了一种新的结构化表达方式。与传统 ANN 和 DNN 相比，GNN 能够将阵元、阵元间距、耦合关系和拓扑连接纳入统一模型中，突破了规则阵列或固定拓扑结构的限制。因此，将图深度学习应用于天线阵列分析、综合和优化具有较好的理论基础和发展潜力。

\subsection{基于元启发式算法的天线设计研究现状}

深度学习模型和机器学习模型在天线设计中多用作代理模型，其作用是替代传统全波仿真方法，对天线性能或阵列方向图进行快速预测；而元启发式算法则更多用于天线或阵列参数的反向寻优。元启发式算法通常通过模拟自然界或群体行为中的搜索机制进行优化，适合求解高度非线性、多变量和多目标问题，因此在天线设计优化中得到广泛应用。

遗传算法（Genetic Algorithm, GA）和粒子群优化算法（Particle Swarm Optimization, PSO）是较早应用于天线设计的典型元启发式算法，并且发展较为成熟。近年来，研究重点逐渐从单独使用传统优化算法转向“代理模型 + 智能优化算法”的协同框架。例如，Liu 等人提出的多边形贴片天线多目标逆向设计方法采用 CNN 作为快速代理预测模型，并结合多岛差分进化算法完成多目标搜索，使任意形状贴片天线能够在较短时间内完成综合优化[10]。

随着天线结构和优化目标日益复杂，更多新型元启发式算法被引入天线代理建模和参数优化中。Huang 等人研究了元启发式智能算法优化神经网络权值和偏置的天线建模方法，引入海鸥优化算法、融合反向学习的改进蝴蝶算法和人工兔优化算法等，对神经网络进行优化。实验结果显示，融合反向学习的改进蝴蝶算法优化后的神经网络在预测精度和运行效率方面表现更优，说明新型群智能算法可用于提升天线代理模型的建模精度[11]。

此外，为克服高保真电磁仿真计算成本高的问题，Koziel 等人在 2023 年研究了面向天线优化的变保真元启发式优化流程，通过在不同分辨率电磁模型之间切换来降低自然启发式优化过程中的计算开销[12]。2025 年的相关研究进一步将 1D GP-HetCNN 代理模型与 WOA、COA、GWO、PSO 和贝叶斯优化等方法结合，用于超宽带天线参数优化；在统一代理模型评估预算下，不同智能算法通过调用训练好的代理模型搜索更优几何参数，仅需少量 HFSS 验证即可完成优化流程[13]。这类研究表明，智能优化算法正在与轻量化代理模型、多保真模型和深度学习模型更紧密地结合。

总体而言，元启发式算法已在天线结构设计、阵列综合、带宽展宽、旁瓣抑制和多目标优化等任务中得到广泛应用。其优势在于不依赖梯度信息，适合处理复杂非线性问题；不足在于计算过程中可能需要大量性能评估，且部分算法存在局部最优和收敛速度不稳定等问题。因此，如何将元启发式算法与机器学习代理模型、深度学习预测模型相结合，以进一步提高优化效率和设计精度，是当前 AI 天线优化设计中的重要研究方向。

\section{讨论与分析}

\subsection{现有技术优势}

\subsubsection{设计效率提升}

AI 辅助天线设计方法最直接的优势体现在设计效率的提升。传统天线设计通常依赖全波电磁仿真软件对结构参数进行反复扫描和迭代优化，单次仿真往往需要较高的计算资源和较长的计算时间。已有综述指出，机器学习与深度学习方法可通过建立天线结构参数与电磁响应之间的代理映射，提高天线设计、优化和选型效率[1-3]。正如第二章文献所述，将深度学习作为代理模型，能够在训练完成后快速预测天线性能，从而显著减少高保真全波仿真的调用次数。

这种效率提升改变了仿真在设计流程中的作用。例如，利用低成本代理建模技术，可以在极少训练样本下快速完成多个频段的透射阵天线设计，Belen 等仅使用 270 个训练样本和 100 个保留样本构建透射阵单元代理模型，并完成 8-14 GHz、22-28 GHz 和 28-36 GHz 多频段设计与实验验证[15]。对于复杂的优化任务，代理模型还可以与多岛差分进化算法、粒子群优化（PSO）或变保真元启发式流程结合，形成“代理模型 + 智能优化算法”的快速搜索框架[10-12]。在这种协同框架下，大量候选结构的初步筛选和优化迭代由代理模型完成，仅需少量全波仿真（如 HFSS）进行验证，从而极大缩短了设计周期。

\subsubsection{非线性建模与多目标优化能力增强}

天线结构参数与电磁性能之间通常呈现强非线性关系。传统经验公式难以准确描述贴片尺寸、馈电位置和阵元间距微小变化所带来的影响。而人工神经网络（ANN）、径向基函数（RBF）网络和卷积神经网络（CNN）展现出了强大的非线性逼近能力[2,9]。研究表明，RBF 神经网络能够有效适应天线阵列的非线性特性，实现稳健的方向性估计[18]。

在多目标优化方面，AI 方法也表现出极强的适应性。现代天线设计需要同时兼顾反射系数、输入阻抗、辐射方向图和带宽等多个指标。通过深度学习代理模型与元启发式算法（如海鸥优化算法、改进蝴蝶算法等）的深度融合，系统能够在几分钟内完成对多边形贴片天线等多目标逆向设计[10]。此外，新型群智能优化算法也能反过来优化神经网络的权值和偏置，进一步提升非线性代理模型的预测精度[11]。

\subsubsection{复杂结构表征能力提升}

随着天线形态向非均匀阵列、流体天线和复杂拓扑结构发展，传统基于规则欧式数据的参数化表达逐渐失效。AI 方法的重要优势在于突破了这一限制，提供了更多维的结构表征方式。对于任意形状的多边形天线或 3D 结构，CNN 与 LSTM 结合的网络能够直接以图像作为输入进行特征提取并自动生成建模代码[16]。

更重要的是，图神经网络（GNN）的引入为复杂天线拓扑提供了天然的表达工具。在天线阵列中，GNN 能够将阵元、阵元间距以及耦合关系作为图的“节点”和“边”纳入统一模型中。例如，在随机稀疏线阵测向或流体天线系统联合优化中，GNN 通过邻居节点聚合有效补偿了缺失阵元信息，在不对阵列几何形状作过多假设的前提下，显著提升了复杂结构的表征能力和阵列综合效果[19,22-23]。

\subsection{现有技术不足}

\subsubsection{数据样本依赖较强}

尽管深度学习代理模型能够显著加快预测过程，但其本质是数据驱动的，性能高度依赖训练样本的数量和质量。在天线设计中，获取高质量数据需要依赖耗时的全波电磁仿真。这意味着 AI 方法在一定程度上是将计算成本从“优化迭代阶段”转移到了“数据集构建阶段”[1,15]。

为了缓解传统手工建模和扫参获取数据效率低下的问题，现有研究甚至需要专门开发基于 CNN-LSTM 的结构自动建模框架来加速数据生成过程[16]。这侧面反映了高昂的数据获取成本依然是制约深度学习在天线设计中广泛应用的核心瓶颈。对于超表面或大型阵列等高维设计空间，若缺乏有效的主动学习、自适应采样或变保真建模策略，庞大的样本需求将使其失去效率优势[5,12]。

\subsubsection{泛化能力与物理可解释性不足}

现有多数 AI 天线设计模型（如基于 ANN 或单一 CNN 的预测器）主要捕捉训练数据中的统计规律，缺乏明确的物理机理支撑。当待预测的天线拓扑、工作频段或材料边界条件超出训练集分布时，模型的泛化能力会急剧下降。微波 CAD 领域的 ANN 综述也指出，数据集质量、训练样本覆盖范围和模型外推能力仍是神经网络代理模型落地时必须面对的问题[2,9]。

同时，深度学习模型往往是一个“黑盒”，其输出结果虽然在数值上能够拟合 S 参数或方向图，但无法保证这些结果严格满足麦克斯韦方程、能量守恒或互易性定理等基础电磁规律。若模型仅追求拟合精度而脱离物理约束，极易在设计空间的边缘区域产生非物理的预测结果，导致逆向设计失效。物理信息神经网络（PINN）的核心思想正是将偏微分方程等物理规律嵌入损失函数或网络训练过程，为解决此类问题提供了重要启发[24-25]。

\subsubsection{工程验证与可制造性考虑不足}

AI 算法在复杂天线参数反向寻优和阵列综合（如非均匀阵列合成）中可能生成理想化的数值解，但这些解并不一定具备良好的可制造性。许多研究仍停留在软件仿真和算法对标阶段，未将加工公差、介质损耗、装配误差和馈电网络复杂度作为约束条件纳入模型[1,3]。

虽然部分研究通过实验验证了代理模型的有效性[15]，但总体而言，面向 AI 生成结构的暗室测试和长期稳定性评估依然缺乏。若优化结果缺少实物验证环节的闭环校正，全波仿真结果与实际工程性能之间的偏差将限制该类方法的落地应用。

\subsection{研究空白}

\subsubsection{标准化数据集与评价体系不足}

目前的 AI 天线设计研究大多依赖研究者自行构建的数据集，不同文献在天线类型、频段（如 8-14 GHz 或 28-36 GHz）、仿真软件、网格精度和评价指标上各不相同[13,15]。这种“各自为战”的局面导致不同深度学习模型或改进后的元启发式算法（如 WOA、COA、GWO）难以进行横向公平比较，也阻碍了优秀算法的工程复用。构建开源的、涵盖多频段多拓扑的标准天线数据集与统一评价指标，是当前领域的迫切需求。

\subsubsection{物理机理与数据驱动的深度融合不充分}

虽然部分研究引入了多保真度模型来平衡计算精度与开销[5,12]，但大多数模型依然停留在纯数据拟合层面。天线设计本质受制于复杂的电磁场分布和边界条件。如何将电磁学先验知识（如等效电路模型、解析公式或电磁偏微分方程）嵌入到神经网络的损失函数或网络结构中，构建“物理信息神经网络（PINN）”，从而提升小样本条件下的预测精度和跨频段泛化能力，是亟待突破的研究空白[24-25]。

\subsubsection{系统级与多物理场协同优化研究不足}

现有的深度学习和图神经网络研究，大多针对单一天线的 S 参数预测或特定阵列的方向图综合（如波达方向估计或流体天线波束形成）[19,22-23]。然而，在实际工程中，天线系统不仅需要满足电磁性能，还需要在复杂的通信系统中进行功率分配联合优化（如二部图建模联合优化），同时承受热、力、风载等环境影响。已有多物理场与机器学习辅助优化研究表明，结构、电磁和环境约束的联合建模正成为工程化设计的重要方向[5,26]。从单一电磁性能优化，向包含电磁-热-力多物理场耦合以及系统级通信指标联合优化的演进，是当前研究尚未全面覆盖的盲区。

\subsection{本章总结}

本章基于文献综述分析了 AI 辅助天线设计的发展现状。总体而言，CNN、LSTM、GNN 等深度学习模型与改进的元启发式算法协同工作，已成功将传统天线设计从单一的“仿真试错”转化为“高效建模、精准预测与多目标优化联合”的智能框架，显著提升了针对非均匀阵列、多边形贴片等复杂结构的表征与设计能力。然而，该领域仍面临数据样本获取成本高、模型缺乏物理可解释性及泛化能力弱等痛点。未来的研究应聚焦于减少数据依赖的轻量化代理模型、融入麦克斯韦物理规律的机制约束网络，以及建立涵盖制造约束和多物理场的全流程智能验证体系，以推动 AI 天线技术从理论仿真向工程落地迈进。

\section{改进与未来展望}

正如第二章所述，深度学习代理模型、图神经网络以及生成式模型为天线设计提供了新的方法路径，也显著提升了性能预测、结构优化和复杂阵列建模的效率。然而，结合第三章分析可以看出，现有方法仍存在数据依赖较强、跨结构泛化能力不足、物理机理融合不充分以及工程闭环验证不足等问题。为进一步推动 AI 辅助天线设计从任务专用模型走向可信、可迁移和可落地的智能设计体系，未来研究可从以下几个方向展开。

\subsection{从任务专用代理模型向天线基础模型演进}

现有高保真代理模型，如传统神经网络、支持向量机等，多针对特定天线结构、特定频段或特定性能指标建立映射关系，本质上仍属于“一任务一模型”的专用建模范式。这类方法在训练数据覆盖范围内通常能够取得较好的预测效果，但当结构类型、频段范围、材料参数或边界条件发生变化时，模型往往需要重新采样和训练，跨任务迁移能力有限。

未来的一个重要方向是构建面向电磁与天线领域的基础模型。该类模型可在大规模、多类型、多模态的电磁数据上进行预训练，学习天线几何结构、材料参数、馈电方式、电磁场分布、S 参数和远场方向图之间的共性表征规律。在此基础上，面对新的天线结构或新的设计任务时，可通过少量样本微调实现快速适配，从而缓解传统代理模型泛化能力不足的问题。与单一代理模型相比，天线基础模型更强调跨结构、跨频段和跨任务的知识迁移能力，是提升 AI 天线设计通用性的重要方向[27]。

\subsection{从图结构建模向全域电磁应用延伸}

图深度学习在天线领域仍具备广阔优化空间与研究前景，后续可扩充阵列样本类型，融入圆阵、圆环阵等多元拓扑结构，构建通用性更强的图神经网络模型；还可重构方向图图结构数据，调整数据输入输出形式，依托图神经网络独立完成天线阵列综合任务。同时可立足电磁领域非欧式数据处理特点，搭建适配各类研究场景的图数据表征方式，将相关技术延伸至超材料结构、超材料天线设计等更多研究方向，进一步拓宽图深度学习在天线与电磁领域的应用范围[28]。

\subsection{从孤立算法组合向智能闭环设计平台发展}

当前 AI 天线优化流程通常由多个相对独立的环节组成，包括结构建模、样本生成、代理模型训练、智能优化、全波仿真验证和结果修正等。这些环节之间仍需要较多人工干预，设计效率和流程自动化水平受到限制。未来，AI 辅助天线设计需要从单一算法研究走向面向工程应用的智能闭环设计平台。

在该平台中，设计者可输入工作频段、增益、带宽、极化方式、尺寸限制和应用场景等需求，系统自动完成结构生成、代理模型预测、优化算法调用、HFSS/CST 仿真验证和结果反馈。大语言模型智能体可用于调度不同工具和算法，例如自动生成仿真脚本、控制电磁仿真软件、分析仿真结果并选择新的采样点。结合主动学习机制，平台可优先仿真最具价值的候选结构，并将仿真或实测结果反馈给代理模型持续更新。由此形成“需求输入—结构生成—代理预测—全波验证—模型更新—工程反馈”的闭环流程。该方向有望提升 AI 天线设计的自动化、可验证性和工程落地能力[29]。

\section{总结}
本文系统梳理了人工智能辅助天线设计的技术演进轨迹，重点剖析了从人工神经网络、图深度学习到元启发式算法的协同优化方案。本综述的主要贡献在于：深刻揭示了当前天线AI设计高度依赖高质量数据、缺乏物理机理约束以及跨任务泛化能力弱的核心瓶颈，并由此论证了天线设计范式向“基础模型（Foundation Models）”跃升的必然性。

未来的 AI 天线设计将彻底告别孤立、碎片化的“一任务一模型”专用范式，加速向着以大规模电磁数据预训练为核心的“天线基础模型”演进以实现跨结构与跨频段的知识迁移，同时推动图深度学习向圆阵、超材料等“全域电磁拓扑”延伸，并依托大语言模型智能体与主动学习机制构建打通“需求-生成-仿真-反馈”的“智能闭环设计平台”，从而全面实现天线设计的通用化、全域化与全流程自动化落地。

\newpage
\section*{参考文献}
\addcontentsline{toc}{section}{参考文献}
\begin{enumerate}[label={\textnormal{[\arabic*]}},leftmargin=2.8em,itemsep=0.35em]
\item Sarker N, Podder P, Mondal M R H, et al. Applications of machine learning and deep learning in antenna design, optimization, and selection: a review[J]. IEEE Access, 2023, 11: 103890-103915.
\item El Misilmani H M, Naous T, Al Khatib S K. A review on the design and optimization of antennas using machine learning algorithms and techniques[J]. International Journal of RF and Microwave Computer-Aided Engineering, 2020, 30(10): e22356.
\item Khan M M, Hossain S, Mozumdar P, et al. A review on machine learning and deep learning for various antenna design applications[J]. Heliyon, 2022, 8(4): e09317.
\item Li L L, Wang L G, Teixeira F L, et al. DeepNIS: Deep neural network for nonlinear electromagnetic inverse scattering[J]. IEEE Transactions on Antennas and Propagation, 2019, 67(3): 1819-1825.
\item Yao H M, Jiang L J. Machine learning based neural network solving methods for the FDTD method[C]//Proceedings of the 2018 IEEE Antennas and Propagation Society International Symposium on Antennas and Propagation \& USNC/URSI National Radio Science Meeting. Boston: IEEE, 2018: 2321-2322.
\item Zhang, Q. J., Gupta, K. C., \& Devabhaktuni, V. K. “Artificial neural networks for RF and microwave design—From theory to practice.” IEEE Transactions on Microwave Theory and Techniques, 2003.
\item LeCun, Y., Bengio, Y., \& Hinton, G. “Deep learning.” Nature, 2015.
\item Shi D, Fang W, Zhang F F, et al. A novel method for intelligent EMC management using a knowledge base[J]. IEEE Transactions on Electromagnetic Compatibility, 2018, 60(6): 1621-1626.
\item Feng F, Na W, Jin J, Zhang J, Zhang W, Zhang Q J. Artificial neural networks for microwave computer-aided design: the state of the art[J]. IEEE Transactions on Microwave Theory and Techniques, 2022, 70(11): 4597-4619.
\item Singh P, Panda S S, Hegde R S. Deep-learning empowered multi-objective antenna design: a polygon patch antenna case study[C]. 2024 National Conference on Communications (NCC), 2024: 1-6. DOI: 10.1109/NCC60321.2024.10486033.
\item Huang J, Nan J, Gao M, Wang Y. Antenna modeling based on meta-heuristic intelligent algorithms and neural networks[J]. Applied Soft Computing, 2024, 159: 111623.
\item Pietrenko-Dabrowska A, Koziel S, Leifsson L. Expedited metaheuristic-based antenna optimization using EM model resolution management[C]. International Conference on Computational Science, 2023.
\item Yang X, Nan J, Wang M. Antenna modeling and optimization based on 1D GP-HetCNN surrogates and intelligent methods[J]. Discover Applied Sciences, 2026, 8(2): 1-22. DOI: 10.1007/s42452-025-07962-7.
\item Mishra, R. K., \& Patnaik, A. “Neural network-based CAD model for the design of square-patch antennas.” IEEE Transactions on Antennas and Propagation, 1998, 46(12): 1890-1891.
\item Belen M A, Caliskan A, Koziel S, Pietrenko-Dabrowska A, Mahouti P. Optimal design of transmitarray antennas via low-cost surrogate modelling[J]. Scientific Reports, 2023, 13: 15044.
\item Wei Z, Zhou Z, Wang P, Ren J, Yin Y, Pedersen G F, Shen M. Fast and automatic 3D modeling of antenna structure using CNN-LSTM network for efficient data generation[EB/OL]. arXiv:2306.15530, 2023.
\item Ayestaran R G, Las-Heras F, Martinez J A. Non uniform-antenna array synthesis using neural networks[J]. Journal of Electromagnetic Waves and Applications, 2007, 21(8): 1001-1011.
\item Mishra S, Yadav R N, Singh R P. Neural estimations of first null beamwidth for broadside and end-fire uniform linear antenna arrays using RBF neural networks[C]//Advances in Intelligent Systems and Computing. Springer, 2015, 343: 445-451. DOI: 10.1007/978-81-322-2268-2\_46.
\item Yang Y, Zhang M, Peng S, Ye M, Zhang Y. Direction-of-arrival estimation for a random sparse linear array based on a graph neural network[J]. Sensors, 2024, 24(1): 91.
\item Wu, Z., Pan, S., Chen, F., Long, G., Zhang, C., \& Philip, S. Y. “A comprehensive survey on graph neural networks.” IEEE Transactions on Neural Networks and Learning Systems, 2021.
\item Yu, B., Yin, H., \& Zhu, Z. “Spatio-temporal graph convolutional networks: A deep learning framework for traffic forecasting.” International Joint Conference on Artificial Intelligence, 2018.
\item He C, Lu Y, Chen W, Ai B, Wong K K, Niyato D. Graph neural network enabled fluid antenna systems: a two-stage approach[EB/OL]. arXiv:2502.03922, 2025.
\item Xie X, Lu Y, Ding Z. Graph neural network enabled pinching antennas[EB/OL]. arXiv:2502.05447, 2025.
\item Raissi M, Perdikaris P, Karniadakis G E. Physics-informed neural networks: a deep learning framework for solving forward and inverse problems involving nonlinear partial differential equations[J]. Journal of Computational Physics, 2019, 378: 686-707.
\item Karniadakis G E, Kevrekidis I G, Lu L, et al. Physics-informed machine learning[J]. Nature Reviews Physics, 2021, 3: 422-440.
\item Han B, Wu Q, Yu C, et al. Low-wind-load broadband dual-polarized antenna and array designs using sequential multiphysics machine-learning-assisted optimization[J]. IEEE Transactions on Antennas and Propagation, 2025, 73(1): 135-148.
\item 《基于深度学习与代理模型的天线优化设计研究进展》\_汤昊臻.
\item 《基于图深度学习的天线阵列优化设计》\_苏倩.
\item 《基于智能优化算法和代理模型的天线优化设计》\_范文敏.
\end{enumerate}

\end{document}
